% Provenance: copied from the design workspace CambridgePQ/recursion.fast, `report/sections/05-brief-review.tex`, on 2026-10-19. Read-only copy; the source is authoritative.
\section{Critical review of the brief}\label{sec:brief-review}

The brief~\cite{recursionfastbrief} fixes the intent of recursion.fast: a competition in which
untrusted agents submit SNARK verifier specifications for arbitrary R1CS, as RISC-V programs with an
honest prover and a two-to-one \emph{batcher}, together with Lean~4 proofs of knowledge soundness and
completeness, scored by a cost function that should reward interesting, general verifier architectures
rather than trivial or overfitted ones. The intent is sound and the precedents it names (ots.golf,
sig.golf, the Proximity Prize; \cref{sec:prior}) show that the format works. Before designing, we
reviewed the brief for ambiguities and missing scope. \Cref{tab:brief-review} lists the twelve points
we found; each is resolved later in the report, and the last column summarises the resolution and
names the decision (blue box) or section that makes it. The four points with the largest
consequences are discussed first.

\paragraph{The objective and the proposed score.} \enquote{As efficient as possible} admits four
readings that pull against each other: final-layer verifier cost, cost of one recursion step, proof
size, and prover time. Lowering the code rate shrinks the first three and raises the fourth; grinding
does the same in the Fiat--Shamir setting. The name of the competition and the brief's
\enquote{ambitious variant} (arithmetise the verifier and measure 2-to-1 recursion) point at the
recursion step, which we adopt as the primary objective (\cref{sec:objectives}). The proposed score,
\enquote{median performance on hash-chain proofs divided by the size of the RISC-V verifier and
batcher programs}, is dimensionally inverted if performance is measured in cycles: dividing a cost by
a size rewards \emph{larger} programs. Any additive or multiplicative size term would also reward
obfuscation and code compression, which no user pays for. Both earlier golf competitions instead cap
program size and do not score it; ots.golf introduced its cap after launch because an uncapped image
is a free lookup table~\cite{otsgolf2026}. We follow them (\cref{sec:cost}).

\paragraph{Succinctness for arbitrary R1CS requires preprocessing.} A verifier that reads the
constraint matrices pays $\Theta(\mathrm{nnz})$ field operations, which is the cost of the trivial NP
verifier. Succinct verification of a \emph{non-uniform} relation requires an indexer that commits to
the circuit once, producing a short verification key and an index oracle that the verifier queries
(holography; \cref{subsec:bg-r1cs}). The brief does not mention an indexer or a verification key. We
add both, cap the key (otherwise it is again a lookup table) and price every access to the index
(\cref{sec:protocol}).

\paragraph{Knowledge soundness does not prevent overfitting.} A verifier specialised to the
benchmark's circuit can remain knowledge-sound, because the extractor may carry the benchmark
witness; overfitting must instead be excluded by a class-uniform cost theorem
(\cref{subsec:bench-hardcode}). Relatedly, Lean extractors have no running-time bound, so knowledge
soundness collapses to soundness (\cref{subsec:sec-extractor}).

\paragraph{The batcher property as stated is completeness only.} A batcher that outputs a fixed
accepting proof, or concatenates, satisfies it; the load-bearing property is knowledge of witnesses
for \emph{both} inputs with a batched object of bounded size (\cref{subsec:rec-critique}). We define
the batcher as an accumulation scheme in the information-theoretic model and state the recursive
deployment, which no random-oracle theorem covers~\cite{BCG24-RelROM}, as a named assumption
(\cref{subsec:rec-cannot}).

\begingroup
\small
\renewcommand{\arraystretch}{1.15}
\begin{longtable}{@{}p{0.03\textwidth}p{0.30\textwidth}p{0.29\textwidth}p{0.30\textwidth}@{}}
\caption{Points raised in the review of the brief and their resolution.}\label{tab:brief-review}\\
\toprule
\# & Point & Why it matters & Resolution \\
\midrule
\endfirsthead
\toprule
\# & Point & Why it matters & Resolution \\
\midrule
\endhead
1 & Objective ambiguous (final verifier, recursion step, proof size, prover) & the four pull in opposite directions & recursion-step cost primary; proof size and prover time as constraints and secondary tracks (\cref{subsec:obj-what}) \\
2 & Score \enquote{performance $\div$ program size} & dimensionally inverted; size terms reward obfuscation & cycles plus a byte charge; hard 1~MiB program cap (COST-1, \cref{sec:cost}) \\
3 & No indexer or verification key & arbitrary R1CS cannot be verified succinctly without preprocessing & holographic model: capped key, priced index oracle (P-1, \cref{subsec:proto-indexer}) \\
4 & Interactive \enquote{proof} vs argument unspecified & succinct interactive proofs for NP need an oracle or commitment model & interactive: public-coin IOP in Lean; BCS and Fiat--Shamir as cited meta-theorems with explicit bounds (P-1, S-4; \cref{subsec:sec-cap}) \\
5 & Honest prover in RISC-V & a proved RISC-V SNARK prover would dwarf sig.golf's 154k-line proof; but prover cost must still be constrained & Lean prover specification for completeness; static prover caps; native prover for record runs (BEN-2, \cref{subsec:arch-badge}) \\
6 & Knowledge soundness does not prevent hard-coding & extractors are existential & class-uniform cost theorem over all R1CS in a class (BEN-1, \cref{subsec:bench-hardcode}) \\
7 & Batcher property is completeness only & concatenation and constant proofs satisfy it & accumulation scheme with knowledge soundness and a size cap; recursion heuristic named (REC-1, REC-2) \\
8 & Security parameters unspecified (field, bits, proven vs conjectured regimes, quantum, zero knowledge) & Lean proofs cannot assume conjectures; several capacity conjectures are now false & binary relation field $\mathrm{GF}(2^{64})$ with challenges in $\mathrm{GF}(2^{192})$, per-round error $\le 2^{-128}$, proven regimes, named hypotheses in a separate track; classical; no ZK (P-3, S-2, S-5) \\
9 & \enquote{Median over millions of runs} & honest verification cost is essentially input-independent; measurements are noisy and gameable & proved worst-case bound, an 18-counter cost vector, over all accepting runs (fast-failing rejections unscored, fuel-capped); measurements as unscored demonstrations (COST-1, BEN-1) \\
10 & Field arithmetic in RV64IM unpriced & the cost model decides which field and which arithmetic wins & contract machine RV64IM+Zbc (carry-less multiply as the only field assist, priced separately); v0.3 per-profile zkVM cost vectors (recursion board) and a latency board; no hash or field syscall on T1 (a field-syscall board T1-F is planned); field fixed by the relation (P-5, COST-1) \\
11 & Toolchain and check-time feasibility & ArkLib and riscv-zkvm pin different Lean versions; real submissions take hours to check & single pinned toolchain with vendored RISC-V model; measured budgets (A-2, \cref{subsec:base-time}) \\
12 & Trusted vs untrusted boundaries & imports with \code{sorry}, Lean-level attacks, harness code & explicit trusted surface, comparator pipeline, axiom whitelist, binding rule (A-1, A-3; \cref{tab:arch-trusted}) \\
\bottomrule
\end{longtable}
\endgroup

\paragraph{Scope added.} Beyond the twelve points, the blueprint covers three topics the brief does
not request but which a deployment needs: Fiat--Shamir statement binding against weak-FS and
generated-R1CS attacks~\cite{BPW12-WeakFS,Fen26-GenR1CS}, and the hash-in-circuit attack
of~\cite{KRS25} as a named deployment assumption (\cref{subsec:sec-fs,subsec:rec-cannot}); the governance of trusted
mathematical hypotheses such as proximity-gap theorems that are published but not yet formalised
(\cref{subsec:sec-hypotheses}); and contract versioning and record revalidation, which both golf
competitions found expensive after launch (\cref{sec:issues}).
